% =========================================================================
% $IX$ Additional qualifiers and withdrawn statements (lane C).
% =========================================================================
\section{Additional qualifiers, and the statements we had to withdraw}\label{sec:qual}

The reproduction-side limitations are listed at the end of \S\ref{sec:nofactor}
(rate-side numbers are constructive upper bounds; the blue curve is not
reproduced; the rank-one optimality of the wall-side rule rests on repeated
descent, not a proof). This section adds the qualifiers the rate side needs, plus
the list of our own statements that did not survive.

\paragraph{Non-negotiable wording.}
\emph{(1)} Rate-side numbers are upper bounds from local descent; the only
lower bound we prove is at a frozen prior, and its frontier-wide extension is
\emph{withdrawn} because the monotonicity it needs is false (a closed-form
rank-one counterexample and $115$ reachable in-cone counterexamples).
\emph{(2)} The design rule is ``best known, certified global \emph{on the floor
side}''; we never write ``the optimal $F$''.
\emph{(3)} Rank and threshold statements are \emph{spatial} (which task
dimensions are worth sensing); no temporal threshold statement is made here.
\emph{(4)} Every exponent statement names its window; $\gamma_\infty=2\ln2/r$ is
asymptotic and its finite-rate form is the inequality
$\gamma_r/\gamma_1>1/r$ for $\Delta I\le8$ only.
\emph{(5)} The two excesses $x$ (cost axis) and $\Delta I$ (rate axis) have
different origins and are never mixed in a single fit.
\emph{(6)} Everything quantitative is for the published four-state example.
\emph{(7)} The slope $2\ln2/r$ is reverse-water-filling arithmetic, not our
finding; what is ours is the domain, the finite-rate inequality, the intercept
\eqref{eq:gap}, and the rank structure.

\paragraph{Statuses are re-read, not remembered.} The withdrawn frontier-wide
bound above is the one place in this note where a claim's \emph{status} (not its
value) changed after the numbers were frozen: an earlier draft carried it as
``verified but not proved'' on the strength of $540$ full-rank perturbations, and
a later rank-one perturbation refuted it. Every status word in this note is
therefore read off the current evidence logs, not off the draft.

\paragraph{Calibre independence.} The floor-side numbers do not depend on how
perfect measurement is emulated: a sensor scale $s=10^8$ and the limiting
projection model agree to $0.0136\%$ worst case over $144$ designs, so the wall
and the ladder intercept are calibre-independent quantities.

\paragraph{A negative result worth keeping.} The estimator of the local exponent
is not a selector: refitting its amplitude absorbs a wrong exponent, and gating a
design on the estimator's own rank-deficiency flag costs $+35.66\%$ in one cell
while leaving $2/24$ candidates in another. It is a diagnostic, not a decision
rule.

\paragraph{What we got wrong (our own record).} We report these because each of
them changed a number or a claim in this note.
\begin{itemize}\itemsep2pt
\item \emph{``The optimal support is the leading eigenspace of $\Theta$, i.e.\ the
real object is water-filling on $\Theta$'s spectrum.''} Our own measurement
refuted it: forcing the support there costs $+5.0\times10^{-2}$ and
$+9.5\times10^{-2}$ bit at rank one (\S\ref{sec:rate}). What survives is the rank
structure and the statement that the support is \emph{not} diagonal in the
task-weight basis.
\item \emph{A closed-form prefactor for the ladder.} Our first derivation
identified a block of $s(R-R_\infty)$ with the whole object; the two differ by a
factor ranging over $0.002$--$0.97$, so the closed form was withdrawn; the
identity that survives is \eqref{eq:gap}.
\item \emph{A KKT argument for $\Delta>0$ at every $D$.} Refuted by our own SDP:
the multiplier on the cone constraint is non-zero, $I(S)$ is \emph{not} convex
(a midpoint test violates convexity by $1.8\times10^{-2}$ bit), and the
family $S=c\Theta$ is stationary but suboptimal by $0.13$--$0.24$ bit. This is
what forced the reduction in \S\ref{sec:rate} to a uniqueness statement.
\item \emph{A vacuous diagnostic.} Our first ``diagonality'' test used principal
angles at full rank, where it is empty by definition; the replacement is the
pair of measures related by the identity in \S\ref{sec:rate}.
\item \emph{Two bookkeeping errors that changed numbers}: a sweep table indexed by
ragged arrays, i.e.\ misaligned rows; and a sensor family whose rate column had
been paired to a \emph{different} design than its cost column, worth $0.0350$ bit
once the pairing is fixed. Both are corrected here; the second one flipped the
sign of a sub-claim.
\item \emph{A solver artefact.} In the large-$s$ regime our information-form fixed
point loses accuracy because its tolerance is relative to a covariance that grows
with $s$; the DARE route is the reliable one. All large-$s$ numbers quoted here
come from the DARE route.
\end{itemize}

\paragraph{What would change the conclusions.} One counterexample to the strict
case of the KKT statement (a rank-$r<r^{\star}(D)$ subspace carrying a lossless
optimal design) would turn the ``infeasibility interval'' picture from a theorem
into a counterexample family; we look for that counterexample directly and report
the search in Appendix~\ref{app:repro}. One plant for which the two legs
\eqref{eq:legs} disagree would break the measured status of
$\gamma_\infty=2\ln2/r$.
